%% kapa: exam=2 round=19 year=2008 period=1 subject=appraisal_practice subject_ko=감정평가실무 seq=2 number=2 points=35
%% source: Q-Net 감정평가사 기출문제 (cst00309) artlSeq=5024177 file=extracted/1008/5024177/2008년 제19회 감정평가사 2차시험 기출문제/1교시.hwp
%% conversion: kordoc 4.24.51 -> tools/gen2.py (auto); review=proofread-sonnet; flags=tikz
\begin{longquestion}{【문제 2】}
토지 소유자 J씨는 C시 D읍 E리 30 번지 토지에 대하여 토지등기부등본을 첨부하여 감정평가사 S씨에게 아래와 같은 조건으로 부동산자문 의뢰를 하였다. 주어진 자료를 활용하여 물음에 답하시오. (35점) \par
\medskip\noindent (1) 2008. 01. 01.을 가격시점으로 하여 토지의 정상가격을 평가하시오. (5점) \par
\medskip\noindent (2) 2008. 01. 01.을 가격시점으로 하여 적산임료 산정을 위한 토지의 기초가격을 평가하시오. (5점) \par
\medskip\noindent (3) 2008. 01. 01.을 가격시점으로 하여 <자료 3>에 주어진 투자조건 등을 만족하는 토지의 투자가격을 결정하시오. (15점) \par
\medskip\noindent (4) 2008. 09. 21.을 가격시점으로 하여 토지의 정상가격을 평가하시오. (5점) \par
\medskip\noindent (5) 2008. 01. 01.을 가격시점으로 한 앞의 세가지 가격을 가치기준(Valuation Bases)에 따라 비교․설명하시오. (5점) \par
<자료 1-1> 사전조사사항 Ⅰ \par
1) 등기부등본 \par
① 토지등기부등본(의뢰시 첨부서류) \par
{\small\centering
\begin{tabular}{|c|c|c|c|}
\hline
\makecell{소재지번} & \makecell{지목} & \makecell{면적} & \makecell{기타사항} \\
\hline
\makecell{C시 D읍 E리 30 번지} & \makecell{임야} & \makecell{630m\textsuperscript{2}} & \makecell{--} \\
\hline
\end{tabular}\par}

② 건물등기부등본 \par
{\small\centering
\begin{tabular}{|c|c|c|}
\hline
\makecell{소재지번 및 건물번호} & \makecell{건물내역} & \makecell{기타사항} \\
\hline
\makecell{C시 D읍 E리 30 번지} & \makecell{목조 함석지붕 창고 단층 36m\textsuperscript{2}} & \makecell{--} \\
\hline
\end{tabular}\par}

2) 토지대장 \par
{\small\centering
\resizebox{\ifdim\width>\linewidth \linewidth\else\width\fi}{!}{%
\begin{tabular}{|c|c|c|c|c|}
\hline
\multirow{2}{*}{\makecell{토지소재}} & \multirow{2}{*}{\makecell{지번}} & \multicolumn{3}{|c|}{\makecell{토지표시}} \\
\cline{3-5}
 &  & \makecell{지목} & \makecell{면적(m\textsuperscript{2})} & \makecell{사 유} \\
\hline
\makecell{C시 D읍 E리} & \makecell{30} & \makecell{전} & \makecell{300} & \makecell{2007년 12월 28일\newline{}분할되어 본번에 - 1을 부함\newline{}2007년 12월 30일\newline{}임야에서 전으로 등록전환} \\
\hline
\makecell{C시 D읍 E리} & \makecell{30-1} & \makecell{임야} & \makecell{330} & \makecell{2007년 12월 28일\newline{}30번지에서 분할} \\
\hline
\end{tabular}}\par}

3) 건축물대장등본 : C시 D읍 E리 30 번지 및 동소 30-1번지는 건물이 등재되어 있지 않음. \par
4) 토지이용계획확인원 : 관리지역, 토지거래계약에 관한 허가구역 \par
5) 지적도 \par
\begin{center}
\begin{tikzpicture}[x=0.0125cm,y=-0.0125cm,font=\scriptsize]
  \draw (324,812) -- (324,958);
  \draw (745,812) -- (745,958);
  \draw (247,958) -- (812,958);
  \draw (247,1007) -- (812,1007);
  \draw (324,1007) -- (324,1105) -- (291,1105) -- (291,1160);
  \draw (535,1007) -- (535,1160);
  \draw (745,1007) -- (745,1160);
  \draw (247,1160) -- (812,1160);
  \draw (369,1160) -- (369,1208);
  \draw (646,1160) -- (646,1208);
  \node at (280,880) {27 답};
  \node at (535,880) {29 장};
  \node at (795,880) {29-1 전};
  \node at (535,983) {569 도};
  \node at (280,1071) {28 답};
  \node at (430,1071) {30-1 임};
  \node at (640,1071) {30 전};
  \node at (785,1071) {36 전};
  \node at (513,1184) {360 전};
\end{tikzpicture}
\end{center}
<자료 1-2> 사전조사사항 Ⅱ \par
1) 인근의 비교가능한 표준지 공시지가 \par
(공시기준일 : 2008. 01. 01.) \par
{\small\centering
\begin{tabular}{|c|c|c|c|c|c|c|c|c|}
\hline
\makecell{연\newline{}번} & \makecell{소재지} & \makecell{면적\newline{}(m\textsuperscript{2})} & \makecell{지목} & \makecell{이용\newline{}상황} & \makecell{용도\newline{}지역} & \makecell{도로\newline{}교통} & \makecell{형상․ 지세} & \makecell{공시지가\newline{}(원/m\textsuperscript{2})} \\
\hline
\makecell{1} & \makecell{C시 D읍 E리 23} & \makecell{455} & \makecell{전} & \makecell{전} & \makecell{관리} & \makecell{세로가} & \makecell{부정형완경사} & \makecell{62,000} \\
\hline
\makecell{2} & \makecell{C시 D읍 E리 50} & \makecell{766} & \makecell{답} & \makecell{답} & \makecell{관리} & \makecell{맹지} & \makecell{부정형\newline{}평지} & \makecell{51,000} \\
\hline
\makecell{3} & \makecell{C시 D읍 E리 135} & \makecell{356} & \makecell{대} & \makecell{단독\newline{}주택} & \makecell{관리} & \makecell{세로가} & \makecell{부정형평지} & \makecell{95,000} \\
\hline
\makecell{4} & \makecell{C시 D읍 E리 150} & \makecell{420} & \makecell{차} & \makecell{주차장} & \makecell{관리} & \makecell{세로가} & \makecell{부정형\newline{}평지} & \makecell{68,000} \\
\hline
\makecell{5} & \makecell{C시 D읍 E리 200} & \makecell{300} & \makecell{대} & \makecell{상업용} & \makecell{관리} & \makecell{소로\newline{}한면} & \makecell{정방형평지} & \makecell{190,000} \\
\hline
\makecell{6} & \makecell{C시 D읍 E리 356} & \makecell{836} & \makecell{임} & \makecell{토지\newline{}임야} & \makecell{관리} & \makecell{세로(불)} & \makecell{부정형\newline{}완경사} & \makecell{43,000} \\
\hline
\makecell{7} & \makecell{C시 D읍 E리 산12} & \makecell{4,260} & \makecell{임} & \makecell{자연림} & \makecell{관리} & \makecell{세로가} & \makecell{부정형\newline{}완경사} & \makecell{30,000} \\
\hline
\end{tabular}\par}

2) 지가변동률 : 국토해양부장관 발표자료로 추정한 2008. 01. 01.부터 2008. 09. 21.까지의 C시 관리지역 지가변동률은 1.000\%임. \par
3) 가격자료 및 기타사항 \par
① 당해지역의 공시지가 수준은 적절한 균형을 유지하고 있으며 적정지가를 비교적 잘 반영하고 있으나 2008. 03. 02. 이후 대지에 대한 수요증가로 국지적인 가격변동이 있었음. \par
② 실거래가는 일부 포착되었으나 세부내역이 없어 검토가 어려움. \par
③ 신뢰할 만한 평가선례자료는 다음과 같음. \par
{\small\centering
\begin{tabular}{|c|c|c|c|c|c|c|c|c|}
\hline
\makecell{연번} & \makecell{소재지} & \makecell{면적\newline{}(m\textsuperscript{2})} & \makecell{지목} & \makecell{이용\newline{}상황} & \makecell{용도\newline{}지역} & \makecell{평가\newline{}목적} & \makecell{가격시점} & \makecell{평가가격\newline{}(원/m\textsuperscript{2})} \\
\hline
\makecell{1} & \makecell{C시 D읍 E리 140} & \makecell{455} & \makecell{대} & \makecell{단독} & \makecell{관리} & \makecell{경매} & \makecell{’08. 07. 11.} & \makecell{140,000} \\
\hline
\makecell{2} & \makecell{C시 D읍 E리 225} & \makecell{766} & \makecell{대} & \makecell{상업용} & \makecell{관리} & \makecell{경매} & \makecell{’08. 06. 20.} & \makecell{300,000} \\
\hline
\end{tabular}\par}

평가선례 및 기타자료 등을 종합검토한 바 2008. 09. 21. 기준 대상토지 평가시 기타요인 보정필요성이 제기되었으며, 분석결과 그 수치는 1.30으로 산정되었음. \par
<자료 2-1> 현장조사사항 Ⅰ : 2008. 01. 01. 기준 \par
1) C시 D읍 E리 30 번지 \par
① 인접토지와 등고평탄한 토지로 현재 지표 위에 부직포를 덮고 주차장부지로 이용중이고 유의할 만한 다른 물건은 없었음. \par
② 당해토지는 북측에 인접한 A공장에 일시적으로 주차장부지로 임대중이라 하며 제시받은 임대차계약서 내용은 아래와 같음. \par
- 소재지 : C시 D읍 E리 30 번지 \par
- 당사자 : 임대인 J, 임차인 A공장 대표이사 R \par
- 임대면적 및 용도 : 300m\textsuperscript{2}, 주차장부지 \par
- 임대금액 : 금---,---원 \par
- 임대기간 : 2008. 01. 01. ～ 2008. 12. 31. \par
- 기타사항 : 임대기간은 J씨의 사정에 의해 임의로 종료될 수 있고 이에 따른 부담은 없으며, 임대 종료시 임차인이 설치한 주차관련 지장물(부직포 등)은 임차인이 제거하기로 함. \par
2) C시 D읍 E리 30-1 번지 \par
① 인접토지와 등고평탄한 부정형토지로 현재 전으로 이용중이고 남서측 일부에는 P씨의 종중묘지가 소재하고 있어 이를 확인한 바 면적은 30m\textsuperscript{2}이고 보존가치가 있어 보존묘지로 지정되어 있는 것으로 조사됨. \par
② 소유자에 따르면 분할 전(前) 30번지는 수년전에 전으로 개간되었고 농지원부에도 등재되어있다 하며 이는 사실로 확인됨. \par
③ 이 토지는 이웃에 거주하는 P씨에게 임대중인 것으로 조사되었으며 제시받은 임대차계약서 내용은 아래와 같음. \par
- 소재지 : C시 D읍 E리 30-1 번지 \par
- 당사자 : 임대인 J, 임차인 P \par
- 임대면적 및 용도 : 330m\textsuperscript{2}, 전 \par
- 임대금액 : 금---,---원 \par
- 임대기간 : 2008. 01. 01. ～ 2008. 12. 31. \par
- 기타사항 : 임대계약은 기간 중 J가 임의로 해지할 수 있고 수년간 임대해 온 점을 고려하여 별도의 부담은 없도록 함. \par
<자료 2-2> 현장조사사항 Ⅱ : 2008. 09. 21. 기준 \par
1) 현장조사시 건축공정률이 80\%정도인 상업용건물을 신축 중이었고 부지조성공사는 완료된 상태였음. \par
2) 제시받은 건축허가서 내용 \par
① 건축구분 : 신축 \par
② 대지위치 : C시 D읍 E리 30, 동소 30-1 \par
③ 대지면적 : 560m\textsuperscript{2} \par
④ 주용도 : 제1종근린생활시설(소매점) \par
⑤ 건축물내역 : 경량철골조 판넬지붕, 1층, 연면적 200m\textsuperscript{2} \par
⑥ 허가번호 : 2008-도시건축과-신축허가-5 \par
⑦ 허가일자 : 2008. 06. 30. \par
⑧ 부속 협의조건 : 종전토지 중 40m\textsuperscript{2}는 도로로 기부체납하고 사업부지는 사업완료 후 지목변경하여야 함. \par
\begin{center}
\begin{tikzpicture}[x=0.0125cm,y=-0.0125cm,font=\scriptsize]
  \fill[gray!45] (324,825) rectangle (745,940);
  \draw (324,622) -- (324,757);
  \draw (745,622) -- (745,757);
  \draw (247,757) -- (812,757);
  \draw (247,802) -- (812,802);
  \draw (324,802) -- (324,888) -- (291,888) -- (291,940);
  \draw (535,802) -- (535,940);
  \draw (745,802) -- (745,940);
  \draw (247,940) -- (812,940);
  \draw (369,940) -- (369,985);
  \draw (646,940) -- (646,985);
  \draw[<->,thin] (273,757) -- (273,802);
  \node at (273,779) [fill=white,inner sep=0.5pt] {4m};
  \draw[<->,thin] (594,802) -- (594,825);
  \node[anchor=west] at (600,813) {1m};
  \draw[->,thin] (310,950) -- (310,925);
  \node at (300,967) {분묘소재부분};
  \node at (280,686) {27 답};
  \node at (535,686) {29 장};
  \node at (795,686) {29-1 전};
  \node at (535,779) {569 도};
  \node at (280,859) {28 답};
  \node at (430,859) {30-1 임};
  \node at (640,859) {30 전};
  \node at (785,859) {36 전};
  \node at (513,962) {360 전};
\end{tikzpicture}
\end{center}
※ 음영부분은 건축허가서상 사업부지임 \par
<자료 3> 투자가격결정을 위한 참고자료 \par
1) 투자조건 \par
① 투자대상 : 숙박시설(모텔 : 객실 30개) \par
② 투자조건 : 가격시점에서 소득수익률≥ 15\%이면 투자 \par
(단, 소득수익률 = 순영업소득(NOI) / 부동산평가액(V)으로 하되, 부동산평가액은 공시지가를 기준한 토지가액과 원가법에 의한 건물가액으로 판단하기로 함.) \par
2) 기타 조사자료 및 참고사항 \par
① 인근의 숙박업소 조사내역 \par
인근의 숙박업소에 대하여 규모를 제외한 가격요인보정 후의 안정화된 소득자료는 다음 표와 같으며 현재까지는 신뢰할만한 것으로 보인다. 감정평가사 S씨의 선임평가사는 가능총소득(PGI), 객실점유율 등의 자료를 분석하여 적절히 활용할 것을 권고하였다. 추세가 있는 경우에는 회귀분석법(Regression Analysis)을 적용하되, 평가서에 세밀한 계산과정은 기술하지 않아도 무방하다고 조언하였다. (단, 구축모형은 유의하다고 가정하고 회귀계수와 회귀상수는 소수점 둘째짜리까지 산정함) \par
{\small\centering
\begin{tabular}{|c|c|c|c|}
\hline
\makecell{조사시점} & \makecell{규모(객실수)} & \makecell{PGI / 객실․월(천원)} & \makecell{객실점유율(\%)} \\
\hline
\makecell{’06. 12.} & \makecell{15} & \makecell{700} & \makecell{75.5} \\
\hline
\makecell{’07. 02.} & \makecell{30} & \makecell{800} & \makecell{82.2} \\
\hline
\makecell{’07. 03.} & \makecell{31} & \makecell{810} & \makecell{81.8} \\
\hline
\makecell{’07. 05.} & \makecell{16} & \makecell{690} & \makecell{74.1} \\
\hline
\makecell{’07. 06.} & \makecell{30} & \makecell{800} & \makecell{80.5} \\
\hline
\makecell{’07. 08.} & \makecell{29} & \makecell{790} & \makecell{80.0} \\
\hline
\makecell{’07. 10.} & \makecell{15} & \makecell{710} & \makecell{72.2} \\
\hline
\makecell{’07. 12.} & \makecell{30} & \makecell{800} & \makecell{78.8} \\
\hline
\end{tabular}\par}

※ 조사시점은 매월 초일을 기준하고, 충당금은 무시함. \par
② 기타소득은 자판기 등의 수익으로 1만원 / 월․객실을 거둘 수 있을 것으로 본다. \par
③ 운영경비(OE)는 제반자료를 분석한 바 아래와 같이 의미있는 결과를 얻을 수 있었다. \par
$\hat{y}$ = 1,200,000원 + 0.4x ($\hat{y}$ : 운영경비, x : 가능총소득, R² = 0.951) \par
\begin{center}
회귀상수($a$) $=\dfrac{\sum y\cdot\sum x^{2}-\sum x\cdot\sum xy}{n\sum x^{2}-(\sum x)^{2}}$ \par\medskip
회귀계수($b$) $=\dfrac{n\sum xy-\sum x\cdot\sum y}{n\sum x^{2}-(\sum x)^{2}}$
\end{center}
④ 환원이율은 신뢰할 만하고 의미있는 시장자료를 분석한 결과 다음과 같이 적용가능한 결과를 얻을 수 있었다. \par
{\small\centering
\begin{tabular}{|c|c|c|c|}
\hline
\makecell{구 분} & \makecell{지가변동성 낮음} & \makecell{지가변동성 중립} & \makecell{지가변동성 높음} \\
\hline
\makecell{토지환원이율} & \makecell{8\%} & \makecell{10\%} & \makecell{12\%} \\
\hline
\makecell{건물환원이율} & \makecell{10\%} & \makecell{11\%} & \makecell{12\%} \\
\hline
\makecell{발생확률} & \makecell{10\%} & \makecell{40\%} & \makecell{50\%} \\
\hline
\end{tabular}\par}

⑤ 건축업자가 제시한 숙박시설의 건물투자비용은 730,000,000원이고 구성항목은 다음과 같으나 건물평가시 원가법에 의한 건물평가항목으로 인정하기 어려운 것은 조정이 필요함. \par
{\small\centering
\begin{tabular}{|c|c|c|c|}
\hline
\multicolumn{2}{|c|}{\makecell{구 분}} & \makecell{비율(\%)} & \makecell{내 역} \\
\hline
\multirow{7}{*}{\makecell{직접\newline{}투자비}} & \makecell{1. 설계비} & \makecell{4.0} & \makecell{감리비, 설계비용} \\
\cline{2-4}
 & \makecell{2. 기본건축비} & \makecell{25.0} & \makecell{기초 및 골조공사비 등} \\
\cline{2-4}
 & \makecell{3. 내외장공사비} & \makecell{25.0} & \makecell{미장, 창호공사 등} \\
\cline{2-4}
 & \makecell{4. 기계설비비} & \makecell{18.0} & \makecell{냉난방, 엘리베이터 등} \\
\cline{2-4}
 & \makecell{5. 전기설비비} & \makecell{10.0} & \makecell{전기 및 통신 공사비 등} \\
\cline{2-4}
 & \makecell{6. 집기비품} & \makecell{4.0} & \makecell{비품, 소모품 등 동산} \\
\cline{2-4}
 & \makecell{소계} & \makecell{86.0} &  \\
\hline
\multirow{3}{*}{\makecell{간접 투자비}} & \makecell{7. 일반관리비} & \makecell{3.0} &  \\
\cline{2-4}
 & \makecell{8. 이윤, 기타} & \makecell{5.0} & \makecell{이윤, 건설이자 등} \\
\cline{2-4}
 & \makecell{소계} & \makecell{8.0} &  \\
\hline
\multirow{3}{*}{\makecell{개업비}} & \makecell{9. 개업준비금} & \makecell{4.0} & \makecell{개업전 인건비, 판촉비 등} \\
\cline{2-4}
 & \makecell{10. 운영자금} & \makecell{2.0} & \makecell{초기운전자금} \\
\cline{2-4}
 & \makecell{소계} & \makecell{6.0} &  \\
\hline
\multicolumn{2}{|c|}{\makecell{총계}} & \makecell{100.0} &  \\
\hline
\end{tabular}\par}

<자료 4> 기타 참고사항 \par
1) 지역요인 : 동일함 \par
2) 개별요인 : 이용상황이 동일하면 별도의 지목감가는 하지 아니함. \par
① 도로접면 \par
{\small\centering
\begin{tabular}{|c|c|c|c|c|}
\hline
\makecell{구분} & \makecell{소로한면} & \makecell{세로가} & \makecell{세로(불)} & \makecell{맹지} \\
\hline
\makecell{소로한면} & \makecell{1.00} & \makecell{0.93} & \makecell{0.86} & \makecell{0.83} \\
\hline
\makecell{세로가} & \makecell{1.07} & \makecell{1.00} & \makecell{0.92} & \makecell{0.89} \\
\hline
\makecell{세로(불)} & \makecell{1.16} & \makecell{1.09} & \makecell{1.00} & \makecell{0.96} \\
\hline
\makecell{맹지} & \makecell{1.20} & \makecell{1.12} & \makecell{1.04} & \makecell{1.00} \\
\hline
\end{tabular}\par}

② 형상 \par
{\small\centering
\begin{tabular}{|c|c|c|c|c|}
\hline
\makecell{구분} & \makecell{정방형} & \makecell{장방형} & \makecell{사다리형} & \makecell{부정형} \\
\hline
\makecell{정방형} & \makecell{1.00} & \makecell{0.99} & \makecell{0.98} & \makecell{0.95} \\
\hline
\makecell{장방형} & \makecell{1.01} & \makecell{1.00} & \makecell{0.99} & \makecell{0.96} \\
\hline
\makecell{사다리형} & \makecell{1.02} & \makecell{1.01} & \makecell{1.00} & \makecell{0.97} \\
\hline
\makecell{부정형} & \makecell{1.05} & \makecell{1.04} & \makecell{1.03} & \makecell{1.00} \\
\hline
\end{tabular}\par}

③ 지세 \par
{\small\centering
\begin{tabular}{|c|c|c|}
\hline
\makecell{구분} & \makecell{평지} & \makecell{완경사} \\
\hline
\makecell{평지} & \makecell{1.00} & \makecell{0.97} \\
\hline
\makecell{완경사} & \makecell{1.03} & \makecell{1.00} \\
\hline
\end{tabular}\par}

\end{longquestion}
